\documentclass{article}
\usepackage[T1]{fontenc}
\usepackage{lmodern}
\usepackage{graphicx}
\usepackage{amsmath,amssymb}
\usepackage[a4paper,margin=2cm]{geometry}
\usepackage[absolute,overlay]{textpos}
\setlength{\TPHorizModule}{1mm}
\setlength{\TPVertModule}{1mm}
\usepackage[indonesian]{babel}
\usepackage{hyperref}
\setlength{\emergencystretch}{2em}
% unit: Halaman 51

\begin{document}

% === Halaman 51 (python/native) ===

Prinsip Confusion dan Diffusion

• Claude Shannon dalam makalah klasiknya tahun 1949, 

Communication theory of secrecy systems, 

memperkenalkan prinsip confusion dan diffusion untuk 

membuat serangan berbasis statistik terhadap cipher 

menjadi lebih sulit dilakukan.

• Dua  prinsip tersebut, confusion dan diffusion, menjadi 

panduan dalam merancang algoritma kriptografi, baik 

stream cipher maupun block cipher.

51

% deskripsi: gambar native xref 211
\noindent\includegraphics[width=0.85\textwidth]{../../images/python/page_051_img_01.jpeg}

\end{document}
